\newcommand{\DomainHeader}[2]{
  \chapter{#1}
  \noindent\textbf{Exam Weight:} #2\par\vspace{0.4em}
}

\newenvironment{Objectives}{\section{Objectives}\begin{itemize}[leftmargin=*]}{\end{itemize}}
\newenvironment{Resources}{\section{Recommended resources}\begin{itemize}[leftmargin=*]}{\end{itemize}}

% Margin tag tying a passage to numbered study-guide objectives, e.g.
% \objref{1.6} or \objref{1.10, 2.5}. Bare numbers; domain prefixes implied
% by the guide's numbering (1=Composition ... 8=Visualization).
% Each number is a link to its row in the Coverage Map (\objrow anchors).
\makeatletter
\newcommand{\obj@strip}[1]{\edef\obj@n{\zap@space#1 \@empty}}
\newcommand{\objref}[1]{\marginnote{%
  \raggedright\scriptsize\sffamily
  \def\obj@sep{}%
  \@for\obj@n:=#1\do{%
    \expandafter\obj@strip\expandafter{\obj@n}%
    \obj@sep\hyperlink{obj:\obj@n}{\textcolor{black!55}{\obj@n}}%
    \def\obj@sep{, }}%
}}
\makeatother
% Coverage-map row anchor: prints the objective number and registers the
% jump target for \objref, raised a line so the row lands below the jump.
\newcommand{\objrow}[1]{%
  \raisebox{1.2\baselineskip}[0pt][0pt]{\hypertarget{obj:#1}{}}#1}

% Connector for the "code -> produces -> scene description" idiom (adopted
% from the official OpenUSD tutorials). Mark authored lines in the second
% listing by starting them with ";;".
\newcommand{\produces}{\par\nopagebreak\noindent\textit{\small produces}\par\nopagebreak\vspace{0.1em}}

% Compact question header for the practice chapters (replaces the roomy
% starred subsections) and a bullet-free, tight answer-choice list — the
% A./B./C. labels make bullets redundant.
\newcommand{\Question}[1]{\par\vspace{0.9em plus 0.3em}%
  \noindent\textbf{Question #1}\par\nopagebreak\vspace{0.2em}\ignorespaces}
\newenvironment{Choices}
  {\begin{itemize}[leftmargin=1.5em, label={}, itemsep=1pt, topsep=2pt,
    parsep=0pt]}
  {\end{itemize}}

\newcommand{\CheatSheetTitle}{\section*{Cheat Sheet (one page)}\addcontentsline{toc}{section}{Cheat Sheet}}
\newcommand{\DrillsTitle}{\section*{Exercises and Drills}\addcontentsline{toc}{section}{Exercises and Drills}}

\newcommand{\AnswerLines}[1]{%
  \par\nopagebreak\vspace{0.2em}\nopagebreak\noindent
  \rule{\textwidth}{0.4pt}\par
  \vspace{0.2em}
  \noindent\textit{Write here:}\par
  \vspace{0.2em}
  \foreach \i in {1,...,#1}{\noindent\rule{\textwidth}{0.4pt}\par\vspace{0.55em}}%
}
